\documentclass{/home/hannesheimberger/Documents/LaTeX/Universität/Vorlesungstemplate_English} %Simon: c:/Users/sirio/Documents/LaTeX/Universität/2_SoSe_2025/Vorlesungstemplate, Hannes: /home/hannesheimberger/Dokumente/LaTeX/Universität/Vorlesungstemplate, Emanuele: /home/stundipper/documents/school/2_SoSe_25/Vorlesungstemplate, Simon widehat:

\setcounter{usewhichtitle}{1} %1 heißt, dass der Vorlesungsname übernommen wird, 0 heißt, dass der gesetzte Titel übernommen wird
\setcounter{booltranslate}{0}
\setcounter{boolbibliography}{0}
\setcounter{boolappendix}{1}
\setcounter{boollecture}{1}
\setcounter{boolalg}{0}
\setcounter{boolmath}{1}
\setcounter{booltoc}{1}
\setcounter{booltitlepage}{1}
\setcounter{boolheads}{1}
\setcounter{boolchapnumbereddefs}{0}
\setcounter{boolthirdauth}{0}
\setcounter{boolgerman}{0}

\addbibresource{Vorlesungsmitschrift.bib}
\renewcommand{\aName}{Hannes Heimberger}
\renewcommand{\instName}{Mathematisches Institut}
\renewcommand{\Semester}{Summer 2026}
\renewcommand{\LectureName}{Number Theory I}
%\renewcommand{\LectureSubtitle}{}
\renewcommand{\ProfName}{Prof. Dr. Shuntaro Yamagishi}
\renewcommand{\groupNum}{\ }
\renewcommand{\titel}{Lecture Notes as of \today}
\renewcommand{\thanks}{I thank Daniel Emse for providing his template, which I adapted for these notes.}

\newcommand{\Nzero}{\mathbb{N}}
\renewcommand{\N}{\mathbb{N}_{\geq 1}}

\begin{document}
\chapter*{Notations and conventions}
\section*{Number sets}
\begin{itemize}
  \item The natural numbers, denoted by $\Nzero$, contain $0$ (which is in accord with \emph{DIN 5473} (\emph{German Industrial Norm}) -- Welcome to Germany! xD).
\end{itemize}
\section*{Subsets and -structures}
\begin{itemize}
  \item If we write $A \subset B$, we mean $A \subseteq B$ and $A \neq B$ (see \emph{DIN 5473}).
  \item We write $\subspace$ or $\subspaceeq$ for substructures (subfields, linear subspaces of vector spaces, etc.) of algebraic structures (didn't yet make it into \emph{DIN 5473} :().
\end{itemize}
For $f: \C \times \cdots \to \C$, we define
\[ \index{$N(f)$}{\color{red} N(f)} = \{ s \in \C: f(s, -) = 0\}\]
and
\[ \index{$N^*(f)$}{\color{red} N^*(f)} = \{ s \in [0,1] \times \R \subset \C: f(s, -) = 0\}.\]
\section*{Quantors}
\begin{itemize}
  \item If we write $\jubelforall_{\dots}$, we mean something holds for almost all $\dots$
\end{itemize}
\section*{Functions}
\begin{itemize}
  \item We write $(x,y) = \gcd(x,y)$.
  \item We have $\{ x \} = x-\lf x\rf \forall_{x \in \R}$.
\end{itemize}
\section*{Fields}
\begin{itemize}
  \item We write $\AA$ for $\overline{\Q}$.
\end{itemize}

\chapter{Geometry of Numbers}
\section{Introduction}
\newLecture{15th April 2026}
\begin{definition}{$\Zprim{2}$}{z_prim_2}
  We denote
  \[ \index{$\Zprim{2}$} {\color{red} \Zprim{2}} = \{ (x,y) \in \Z^2 : (x,y) = 1\}.\]
\end{definition}
\begin{theorem}{Dirichlet, 1842}{dirichlet_1842}
  Let $\alpha \in \R$. Then, for every integer $Q \geq 2$, there exists $(x,y) \in \Z^2_{\text{prim}}$ such that $0 < y \leq Q$ and $|x-\alpha y| \leq \nicefrac{1}{Q}$.
\end{theorem}
\begin{pf}
  Consider the $Q+1$ numbers $0, \alpha-\lf\alpha\rf, 2\alpha-\lf 2\alpha\rf, \dots, Q\alpha-\lf Q\alpha\rf$. Each of these is in the interval $[0,1)$. Partition the interval $[0,1)$ into $Q$ subintervals of length $\nicefrac{1}{Q}$. Now, we can apply the pigeonhole principle to deduce that
  \[\exists_{0 \leq i \leq j \leq Q}: |(j\alpha-\lf j\alpha\rf)-(i\alpha-\lf i\alpha\rf)| \leq \frac{1}{Q}.\]
  This implies
  \[ |(j-i)\alpha - (\lf j\alpha\rf -\lf i\alpha\rf)| \leq \frac{1}{Q}.\]
  Let $\widetilde{y} = j-i$ and $\widetilde{x} = \lf j\alpha\rf - \lf i\alpha \rf$. Then $|\widetilde{y}\alpha-\widetilde{x}| \leq \nicefrac{1}{Q}$. Let $d = (\widetilde{y}, \widetilde{x})$, $y = \nicefrac{\widetilde{y}}{d}$ and $x = \nicefrac{\widetilde{x}}{d}$ and divide through by d to get that
  \[ |y\alpha-x| \leq \frac{1}{dQ} \leq \frac{1}{Q}. \qedhere\]
\end{pf}
A very simple consequence of this:
\begin{corollary}{Irrational approximation}{irr_app}
  Let $\alpha \in \R \setminus \Q$. Then there exist infinitely many $(x,y) \in \Zprim{2}$ such that $y > 0$ and
  \[ \left|\alpha-\frac{x}{y}\right| \leq \frac{1}{y^2}.\]
\end{corollary}
\begin{pf}
  For each $Q \geq 2$, let $(x_Q, y_Q) \in \Zprim{2}$ be such that $0 < y_Q \leq Q$ and $|x_Q-y_Q\alpha| \leq \nicefrac{1}{Q}$, which we know to exist from Theorem \ref{thm:dirichlet_1842}. This implies
  \[ \left| \alpha-\frac{x_Q}{y_Q}\right| \leq \frac{1}{Qy_Q} \leq \frac{1}{y_Q^2}.\]
  Suppose the set
  \[ S = \{ (x_Q,y_Q): Q \geq 2\}\]
  is finite. Then there exists $(x_0,y_0) \in two$ such that $(x_Q,y_Q) = (x_0,y_0) \jubelforall_{Q \geq Q_0}$ for some $Q_0$. This means
  \[ \left| \alpha-\frac{x_Q}{y_Q}\right| = \left| \alpha - \frac{x_0}{y_0}\right| \leq \frac{1}{Q} \jubelforall_{Q \geq Q_0}.\]
  Taking the limit $Q \to \infty$ gives
  \[ \alpha = \frac{x_0}{y_0},\]
  which is a contradiction due to the fact that $\alpha$ was supposed to be irrational. \Lightning
\end{pf}
\begin{question}{Is this error term optimal?}{is_opt_irr_app}
  Can we do better than $\nicefrac{1}{y^2}$ for some $\alpha \in \R^2$? I.e. do there exist infinitely many $(x,y) \in \Zprim{2}$ such that $y > 0$ and
  \[ \exists_{\delta > 0}:\left| \alpha - \frac{x}{y}\right| < \frac{1}{y^{2+\delta}}.\]
\end{question}
\begin{answertoquestion} Yes, but it's rare. However, this is not the case if $\alpha$ is algebraic over $\Q$.
\end{answertoquestion}
  \begin{theorem}{Roth, 1955}{roth_1955}
    Let $\alpha \in \AA$ and $\delta > 0$. Then
    \[ \#\{ (x,y) \in \Zprim{2}: y > 0, |\alpha-\nicefrac{x}{y}| \leq \nicefrac{1}{y^{2+\delta}}\} < \aleph_0.\]
  \end{theorem}
  Proving this is one of the main goals of this course. Another important result we will prove is Thue's theorem.
  \begin{theorem}{Thue, 1909}{thue_1909}
    Let $F \in \Z[X,Y]$ be a binary form (hence homogenious) of degree $d$. Suppose $\deg F(X,1) = d$ and that $F(X,1)$ has at least three distinct zeroes in $\C$. Let $m \in \Z \setminus \{ 0\}$. Then
    \[ \#\{ (x,y) \in \Z^2: F(x,y) = m\} < \aleph_0.\]
  \end{theorem}
  Also, time permitted, Professor Yamagishi'd like to go through the details of Stanley Yao Xiao's proof of Büchi's problem.\interfootnotelinepenalty=10000 \footnote{Professor Yamagishi: This is a paper by a friend of mine. Actually, it has nothing to do with geometry of numbers. But I want to do it anyway.}
  \begin{problem}{Büchi}{büchi} Let $0 < x_1 < \dots < x_5$ be integers such that
  \[ x_{i+1}^2-2x_i^2 + x_{i-1}^2 = 2 \forall_{i \in \{ 2, \dots, 4\}}.\]
  Then there exists $x_0 \in \Z$ such that $x_i = x_0+i \forall_{1 \leq i \leq 5}$.
  \end{problem}
  This leads to the following refinement of Hilbert's tenth problem:
  There is no algorithm which decides the existence of integer solutions to an arbitrary system of diagonal quadratic equations:
  \[ A\begin{pmatrix}
    x_1^2 \\
    \vdots \\
    x_n^2
  \end{pmatrix} = b.\]
 Professor Yamagishi recommends the book \emp{An introduction to the geometry of numbers} by Cassels. He also recommend a course on \emp{Prime Numbers} taught next semester (Selected Topics in Number Theory). Tutorial times are \emp{Monday at 14:00} and \emp{Tuesday at 12:00}. You can register for these \emp{Tomorrow at 16:00}.
 \newLecture{17th April 2026}
 \section{Lattices and discrete subgroups}
 Let $n \geq 2$. Given $x_1, \dots, x_n \in \R^n$, we say that they are linearly independent over $\R$, if $c_1x_1 + \dots + c_rx_r \neq 0 \forall_{(c_1, \dots, c_r)^t \in \R^r \setminus \{ 0\}}$. We consider the usual topology on $\R^n$ generated by the Euclidean metric. For subsets of $\R^n$, we consider the subspace topology induced by this: We say a subset $\Lambda \subseteq \R^n$ is \emp{discrete}\index{Discrete Subset of $\R^n$} if $\#\Lambda \cap B < \aleph_0$ for every bounded subset $B \subseteq \R^n$. Furthermore, $\Lambda$ is a \define[Discrete Subgroup of $\R^n$]{discrete subgroup} if it is discrete and $c_1x_1 + c_2x_2 \in \Lambda \forall_{c_1,c_2 \in \Z, x_1,x_2 \in \Lambda}$. In this case, we define \index{Rank of a Discrete Subgroup of $\R^n$}${\color{red} \rank(\Lambda)} = r$ as the maximal number $r$ such that $\Lambda$ contains $r$ linearly independent vectors over $\R$. A \define[Basis of a Discrete Subgroup of $\R^n$]{basis} of $\Lambda$ is a set $\{ v_1, \dots, v_r\}$ of linearly independent vectors over $\R$ such that $\Lambda = \{ c_1v_1 + \dots + c_rv_r: c_1, \dots, c_r \in \Z\}$.
 \begin{lemma}{Lattice Basis Construction}{lat_base_const}
   Let $\Lambda$ be a discrete subgroup of $\R^n$ and $\rank(\Lambda) = r \geq 1$. Also assume that $\{ w_1, \dots, w_r\} \subseteq \Lambda$ is linearly independent over $\R$. Then there exists a basis $\{ v_1, \dots, v_r\}$ of $\Lambda$ such that for all $1 \leq k \leq r$, we have
   \[ v_k = \alpha_{k,1}w_1 + \dots + \alpha_{k,k}w_k \forall_{\alpha_{k,1}, \dots, \alpha_{k, k} \in [0,1] \text{ and } \alpha_{k,k} \neq 0}.\]
 \end{lemma}
 \begin{pf}
   Let $1 \leq k \leq r$. We define
   \[ S_k = \{ \alpha_1w_1 + \dots + \alpha_kw_k: \alpha_1, \dots, \alpha_k \in [0,1], \alpha_k \neq 0\}.\]
   Since $\Lambda$ is discrete and $S_k$ is bounded,
   \[ \underbrace{0 <}_{w_k \in \Lambda \cap S_k} \#(\Lambda \cap S_k) \ll 1.\]
   We choose $v_k$ to be a vector from $\Lambda \cap S_k$ with the minimal $\alpha_k$. Let us write
   \[ v_k = \alpha_{k,1}w_1 + \dots + \alpha_{k,k}w_k.\]
   Then for any
    \[ \alpha_1w_1 + \dots + \alpha_kw_k \in \Lambda \cap S_k,\]
    we have $0 < \alpha_{k,k} \leq \alpha_k \leq 1$. It is clear that $\{ v_1, \dots, v_r\}$ is linearly independent over $\R$. For each $0 \leq k \leq r$, let
    \[ M_k = \Lambda \cap \left\{ \sum_{j=1}^k \alpha_jw_j: \alpha_1, \dots, \alpha_k \in \R\right\} \text{ with } M_0 = \{ 0\}.\]
    Note that $M_r = \Lambda$. We now use induction to show that $M_k = \{ c_1v_1 + \dots + c_kv_k: c_1, \dots, c_k \in \Z\} \forall_{0 \leq k \leq r}$. When $k = 0$, the statement is trivially true. Suppose the statement holds for $k-1$. \\
    Let
    \[ x = \sum_{j=1}^k \beta_jw_j \in M_k \subseteq \Lambda \text{ and set } z_k = \lf \frac{\beta_k}{\alpha_{k,k}}\rf \in \Z.\]
    Then, $0 \leq \beta_k - z_k\alpha_{k,k} < \alpha_{k,k}$. Consider $x - z_k\underbrace{v_k}_{\in \Lambda} \in \Lambda$. Also, \begin{align*}
    & x-z_kv_k = \sum_{j=1}^k (\beta_j-z_k\alpha_{k,j})w_j \\
    = \quad & \underbrace{\sum_{j=1}^k \lf \beta_j-z_k\alpha_{k,j}\rf w_j}_{\in \Lambda} + u,\end{align*}
    where
    \[ u = \sum_{j=1}^{k-1} \{\beta_j-z_k\alpha_{k,j}\}w_j+(\beta_k-z_k\alpha_{k,k}w_k).\]
    It is clear that $u \in \Lambda$. Also, if $\beta_k-z_k\alpha_{k,k} > 0$, then $u \in S_k$ (because $\{ \beta_j-z_k\alpha_{k,j}\} \in [0,1]$ and $0 < \beta_k-z_k\alpha_{k,k} < \alpha_{k,k} \leq 1$), which implies $u \in \Lambda \cap S_k$. But this contradicts the minimality of $\alpha_{k,k}$ \Lightning. Hence, $\beta_k = z_k\alpha_{k,k}$. Then,
    \[ x-z_kv_k = \sum_{j=1}^{k-1} (\beta_j-z_k\alpha_{k,j})w_j \in M_{k-1}.\]
    This concludes the induction.
 \end{pf}
 \begin{definition}{(Full) Lattice}{lat}
   A \define[(Full) Lattice]{(full) lattice} in $\R^n$ is a discrete sugroup of $\R^n$ which has rank $n$.
 \end{definition}
 \textbf{Remark:} For our intents and purposes, a \emp{lattice} is always \emp{full}.
 \titremark{Bases of Lattices} If $\Lambda$ is a lattice in $\R^n$, then by the previous lemma, there exists a basis $\{ v_1, \dots, v_n\}$ such that $\Lambda = \{ c_1v_1 + \dots + c_nv_n: c_1, \dots, c_n \in \Z\}$.
 \begin{definition}{Determinant of a Lattice}{det_lat}
   We define the \define[Determinant of a Lattice]{determinant} of $\Lambda$ to be
   \[ \det(\Lambda) = \left| \det(v_1 \cdots v_n)\right|.\]
 \end{definition}
 Recall that $\GL_n(\Z)$ is the multiplicative group of $n \times n$-matrices over $\Z$ with determinant $\pm 1$.
 \begin{lemma}{Lattice Basis Transformation}{lat_base_trans}
   Let $\Lambda \subseteq \R^n$ be a lattice. Suppose that $\{ v_1, \dots, v_n\}$ and $\{ w_1, \dots, w_n\}$ are two bases of $\Lambda$. Then there exists $u = (u_{i,j})_{i,j \in \{ 1, \dots, n\}} \in \GL_n(\Z)$ such that
   \[ w_i = \sum_{j=1}^n u_{i,j}v_j \forall_{1 \leq i \leq n}.\]
   In particular,
   \[ |\det(w_1 \cdots w_n)| = |\det(v_1 \cdots v_n)|.\]
 \end{lemma}
 \begin{pf}
   Let us define $U = (u_{i,j})_{i,j \in \{ 1, \dots, n\}}$ by
   \[ w_i = \sum_{j=1}^n u_{i,j}v_j \forall_{1 \leq i \leq n}.\]
   Since $w_i \in \Lambda$ and $\{ v_1, \dots, v_n\}$ is a basis of $\Lambda$, we get $u_{i,j} \in \Z \forall_{1 \leq i,j \leq n}$, which ensures $\det(U)\in \Z$. By the same argument, we define $\widetilde{U} = (\widetilde{u}_{i,j})_{i,j \in \{ 1, \dots, n\}}$ analogously. Note that $U^{-1} = \widetilde{U}$. Hence, $\det(\widetilde{U}) = \det(U) = \pm1$.
 \end{pf}
 Let $M,L$ be two lattices in $\R^n$ with $M \subseteq L$. Further let $\{ v_1, \dots, v_n\}$ be a basis of $L$ and $\{ w_1, \dots, w_m\}$ be a basis of $M$. Define $U_0$ as the $n\times n$ matrix that fulfils
 \[ (w_1 \cdots w_n) = (v_1 \cdots v_n)U_0.\]
 Then, we have $U_0 \in M_n(\Z)$. We define the \define[Index of a Sublattice in a Lattice]{index} of $M$ in $L$ by
 \[ [L:M] = |\det U_0| = \left|\frac{\det M}{\det L}\right|.\]
 \emp{Note:} This does not depend on the choice of bases of $M$ and $L$. \\
\begin{example} Let
 \[ \Lambda_0 = \Span_\Z\left\{ \binom{1}{0}, \binom{0}{1}\right\} \text{ and } \Lambda = \Span_\Z\left\{ \binom{2}{0}, \binom{0}{2}\right\}. \qedhere\]
\end{example}
\newLecture{22nd April 2026}
\section{Convex Bodies}
\begin{definition}{Convex Subset of $\R^n$}{conv_subset_rn}
  A set $C \subseteq \R^n$ is \define[Convex Subset of $\R^n$]{convex} if given any $x,y \in C$, we have
  \[ \{ tx+(1-t)y: t \in [0,1]\} \subseteq C.\]
\end{definition}
\begin{definition}{Central Symmetric Convex Body in $\R^n$}{cent_symm_conv_bod_rn}
  A \define[Central Symmetric Convex Body in $\R^n$]{central symmetric convex body} in $\R^n$ is a closed bounded convex $C \subseteq \R^n$ fulfilling $0 \in C^\circ$ and $x \in C \iff -x\in C$.
\end{definition}
\begin{example}
  Consider any norm on $\R^n$. Then the set
  \[ B_1(0) = \{ x \in \R^n: \|x\| \leq 1\}\]
  is a central symmetric convex body.
\end{example}
\noindent Let $C \subseteq \R^n$ be a central symmetric convex body. Define
\[ \index{$C$-norm $\|x\|_C$}{\color{red} \|x\|_C} = \min\{ \lambda \in \R_{\geq 0}: x \in \lambda C\},\]
where $\lambda C = \{ \lambda x: x \in C\}$.
\begin{lemma}{Norm properties of $\|\cdot\|_C$}{norm_props_c_n}
  The map $\|\cdot\|_C$ is
  \begin{enumerate}[(i)]
    \item well-defined,
    \item a norm on $\R^n$ and
    \item $\lambda C = \{ x \in \R^n: \|x\|_C \leq \lambda\}$.
  \end{enumerate}
\end{lemma}
\begin{pf}
  We'll prove (i). Proving the other two statements is left to the reader as an exercise. Clearly, $\|0\|_C = 0$. Let $x \in \R^n \setminus \{ 0\}$. Consider the set
  \[ S = \{ \lambda \in \R_{\geq 0}: x \in \lambda C\}.\]
  Since $0$ is in the interior of $C$, there exists $\delta > 0$ such that
  \[ \{ z \in \R^n: \|z\|_2 \leq \delta\} \subseteq C.\]
  Hence, we have
  \[ x \in \frac{\|x\|_2}{\delta}C \text{, which implies } S \neq \varnothing.\]
  Let $0 \leq \mu = \inf_{\lambda \in S} \lambda$. By the definition of the infimum, we have
  \[ x \in \left(\mu + \frac{1}{m}\right)C \forall_{m \in \N}.\]
  Thus, we have
  \[ \left(\mu + \frac{1}{m}\right)^{-1}x \in C.\]
  Since $C$ is bounded, it must be that $\mu > 0$. Furthermore, since $C$ is closed and
  \[ \left(\mu+\frac{1}{m}\right)^{-1}x \in C \forall_{m \in \N},\]
  we have
  \[ \lim_{m \to \infty}\left(\mu+\frac{1}{m}\right)^{-1}x = \frac{1}{\mu}x \in C,\]
  which implies $x \in \mu C$, i.e. $\mu \in S$.
  \end{pf}
  \section{Minkowski's First Theorem}
  Given a measurable $S \subseteq \R^n$, we define the $n$-dimensional volume $\vol(S) \in \overline{\R_{\geq 0}}$ as $\lambda^n(S)$, the $n$-dimensional Lebesgue measure of $S$. \\
  \textbf{Some key properties:}
  \begin{itemize}
    \item All open and closed sets are measurable.
    \item Bounded measurable sets have finite volume.
    \item For measurable $S$,
    \[ \vol(S) = \int_S 1 \dl{x_1}\cdots\dl{x_n},\]
    if this integral (the $n$-fold Riemann integral) is defined.
    \item The set of all measurable sets in $\R^n$ forms a $\sigma$-algebra over $\R$.
     \item If $\phi: \R^n \to \R^n$ is a linear transformation and $S \subseteq \R^n$ is measurable, then $\phi(S)$ is measurable with
     \[ \vol(\phi(S)) = |\det\phi|\vol(S).\]
     \item The function $\vol$ is $\sigma$-additive.
  \end{itemize}
  Let $\Lambda$ be a full lattice inside of $\R^n$ and $\{ v_1, \dots, v_n\}$ be a basis. We call
  \[ F = \{ \alpha_1v_1+ \dots + \alpha_nv_n: 0 \leq \alpha_1, \dots, \alpha_n < 1\}\]
   a \define[Fundamental Parallel Piped of $\Lambda$]{fundamental parallel} piped of $\Lambda$ (or a \define[Fundamental Domain]{fundamental domain}). Then
  \[ \R^n = \bigcup_{u \in \Lambda} u+F,\]
  where the $u+F$ are pairwise disjoint. Also,
  \[ \vol(F) = \int_{[0,1)^n} |\det[v_1 \cdots v_n]|\dl{\alpha_1}\cdots\dl{\alpha_n} = \det\Lambda.\]
  \begin{theorem}{Minkowski's First Convex Body Theorem, 1896}{mink_first_conv_bod_thm}
    Let $C$ be a central symmetric convex body in $\R^n$ and $\Lambda$ a lattice. Suppose that
    \[ \vol(C) \geq 2^n\det(\Lambda).\]
    Then $C \cap (\Lambda \setminus \{ 0\}) \neq \varnothing$.
  \end{theorem}
  \begin{pf}
    First assume that $\vol(C) > 2^n\det(\Lambda)$. Then the set $\nicefrac{1}{2}C$ satisfies
    \[ \vol\left(\frac{1}{2}C\right) = \frac{1}{2^n}\vol(C) > \det\Lambda.\]
    Let $F$ be a fundamental parallel piped of $\Lambda$ and
    \[ S_a = \frac{1}{2}C \cap (a+F) \forall_{a \in \Lambda}.\]
    Then,
    \[ \sum_{a \in \Lambda}\vol(S_a) = \vol\left(\frac{1}{2}C\right) > \det\Lambda.\]
    Let us define
    \begin{align*} T_a = \quad & -a+S_a \\
    = \quad & \left(-a+\frac{1}{2}C\right) \cap F \subseteq F.
    \end{align*}
    Now,
    \[ \sum_{a \in \Lambda}\vol(T_a) = \sum_{a \in \Lambda}\vol(S_a) > \det\Lambda = \vol(F).\]
    We proved the existence of a collection of subsets of $F$ whose sum of volumes is greater than $\vol(F)$. Therefore, there exists $a \neq b \in \Lambda$ such that $T_a \cap T_b \neq \varnothing$. Let $u \in T_a \cap T_b$. There exists $x,y \in \nicefrac{1}{2}C$ such that $x-a=u=y-b$. \\
    This implies that $x-y = a-b \in \Lambda \setminus \{ 0\}$. We also know that $2x,2y \in C$. Hence, we have $-2y \in C$, which implies
    \[ \frac{1}{2}(2x)+\frac{1}{2}(-2y) = x-y \in C.\]
    Next, let us suppose $\vol(C) = 2^n\det(\Lambda)$. Let $m \in \N$. Since
    \[ \vol\left(\left(1+\frac{1}{m}\right)C\right) > 2^n\det(\Lambda),\]
    it follows from the previous case that there exists
    \[ 0 \neq x_m \in (1+\nicefrac{1}{m})C \cap \Lambda \subseteq 2C\cap\Lambda \forall_{m \in \N}.\]
    We know that $2C$ is bounded, so $2C \cap \Lambda$ is finite. Therefore, there exists an infinite sequence $m_1 < \cdots$ such that $0 \neq x_0 = x_{m_1} = x_{m_2} = \dots$. Since $x_0 = x_{m_i} \in ((1+\nicefrac{1}{m_i})C) \cap \Lambda$, we have
    \[ \left(1+\frac{1}{m_i}\right)^{-1}x_0 \in C \forall_{i \in \N}.\]
    Hence, $0 \neq x_0 \in \Lambda \cap C$.
  \end{pf}
  \begin{lemma}{Cardinality Formula for Lattice-Body Intersections}{card_form_lat_bod_int}
    Let $C$ be a central symmetric convex body and $\Lambda$ a full lattice in $\R^n$. Then
    \[ \#(xC \cap \Lambda) = \frac{\vol(C)}{\det(\Lambda)}x^n + \O(x^{n-1}) \forall_{x \geq 1}.\]
  \end{lemma}
  \begin{pf}
    Let $N(x) = \#(xC\cap \Lambda)$. Consider
    \[ S = \bigcup_{u \in xC \cap \Lambda} u+F.\]
    Since this is a disjoint union, we have
    \[ \vol(S) = N(x)\vol(F).\]
    As $F$ is bounded,
    \[ \exists_{\sigma > 0}: \|y\|_C \leq \sigma \forall_{y \in F}.\]
    Let $x > \sigma$.
    \begin{claim}{}{}
      We have
      \[ (x-\sigma)C \subseteq S \subseteq (x+\sigma)C.\]
    \end{claim}
    The result follows from this claim by taking the volume
    \[ (x-\sigma)^n\vol(C) \leq \vol(S) \leq (x+\sigma)^n\vol(C),\]
    which can be rewritten as
    \[ \vol(S) = x^n\vol(C)+\O(x^{n-1}).\]
    Then just divide by $\vol(F) = \det\Lambda$ (see above).
    \begin{pf}[Proof of the claim]
      Given any $\widehat{x} \in \R^n$, let us write $\widehat{x} = u + y$ with $u \in \Lambda$ and $y \in F$. Then $(x-\sigma)C \subseteq S$, because if $\|\widehat{x}\|_C \leq x-\sigma$, then this implies $\|u\|_C \leq x$ (use the triangle inequality $\|u\|_C \leq \|\widehat{x}-u\|_C + \|\widehat{x}\|_C = \|y\|_C+\|\widehat{x}\|_C \leq \sigma + x-\sigma = x$). Similarly, $S \subseteq (x+\sigma)C$ follows from $\|u\|_C \leq x$, which implies $\|\widehat{x}\|_C \leq x+\sigma$ (again triangle inequality).
    \end{pf}
    This completes the proof.
  \end{pf}
  \newLecture{24th April 2026}
  We will study another proof of Minkowski's first theorem \ref{thm:mink_first_conv_bod_thm} using Lemma \ref{lem:card_form_lat_bod_int}.
  \begin{pf}[Second proof of Theorem \ref{thm:mink_first_conv_bod_thm}]
    Let $m \in \N$. We divide $\Lambda$ into congruence classes modulo $2m$. Each $x \in \Lambda$ is written as
    \[ x=x_1v_1 + \dots + x_nv_n,\]
    where $\{ v_1, \dots, v_n\}$ is a basis of $\Lambda$. Consider
    \[ \Lambda_a = \{ x_1v_1 + \dots + x_nv_n: x_i \equiv a_i \mod 2m (1 \leq i \leq n)\},\]
    for each $a \in \{ 0, \dots, 2m-1\}^n$. Using Lemma \ref{lem:card_form_lat_bod_int}, we have
    \begin{align*}
      \#(mC \cap \Lambda) = \quad & \frac{\vol(C)}{\det(\Lambda)}m^n + \O(m^{n-1}) \\
      > \quad & 2^n(1+\epsilon)m^n,
    \end{align*}
    provided that $m$ is sufficiently large. Hence, there exists $a \in \{ 0, \dots, 2m-1\}^n$ such that $\Lambda_a \cap mC$ contains at least two distinct vectors $u,v$. Then $0 \neq \nicefrac{1}{2m}(u-v) \in \Lambda$. Since we have $u,v \in mC$, we also have $u,-v \in mC$, which implies $\nicefrac{1}{2}u-\nicefrac{1}{2}v \in mC$. Finally, $\nicefrac{1}{2m}(u-v) \in C$, which completes the proof.
  \end{pf}
  \begin{corollary}{}{1.4.4}
    Let $\alpha_{i,j} \in \R$ for $1 \leq i,j \leq n$ and $L_i(x) = \alpha_{i,1}x_1 + \dots + \alpha_{i,n}x_n$. Suppose that $0 < |\det[\alpha_{i,j}]| \leq A_1\cdots A_n$, where $A_1, \dots, A_n > 0$. \\
    Then, there exists $x \in \Z^n \setminus \{ 0\}$ such that
    \[ |L_i(x)| \leq A_i \forall_{i \in \{ 1, \dots, n\}}.\]
  \end{corollary}
  \begin{pf}
    Let $\mathfrak{B} = [-1,1]^n$ and $\Lambda = \left\{ \left(\frac{L_1(x)}{A_1}, \dots, \frac{L_n(x)}{A_n}\right): x \in \Z^n\right\}$. Then this is a lattice inside $\R^n$. Also,
    \[ \Lambda = \left\{ \begin{pmatrix}
      \frac{\alpha_{1,1}}{A_1} & \cdots & \frac{\alpha_{1,n}}{A_1} \\
      \vdots & & \vdots \\
      \frac{\alpha_{n,1}}{A_n} & \cdots & \frac{\alpha_{n,n}}{A_n}
    \end{pmatrix}\begin{pmatrix}
      x_1 \\
      \vdots \\
      x_n
    \end{pmatrix}: x \in \Z^n \right\}.\]
    We have
    \[ \det\Lambda = \frac{|\det[\alpha_{i,j}]|}{A_1\cdots A_n} \leq 1,\]
    which implies
    \[ \frac{\vol(\mathfrak{B})}{\det(\Lambda)} = \frac{2^n}{\det\Lambda} \geq 2^n.\]
    By Minkowski's first theorem, there exists $0 \neq x \in \Lambda \cap \mathfrak{B}$.
  \end{pf}
  \begin{corollary}{}{1.4.5}
    Let $\alpha_1, \dots, \alpha_n \in \R$, not all in $\Q$. Then,
    \[ \# \left\{ (x_1, \dots, x_n,y) \in \Zprim{n+1}: y >0, \left| \alpha_i-\frac{x_i}{y}\right| \leq y^{-1-\nicefrac{1}{n}} \right\} = \aleph_0 \forall_{1 \leq i \leq n}.\]
    (Denote $\Zprim{n} = \{ (x_1, \dots, x_n) \in \Z^n: \gcd(x_1, \dots, x_n) = 1\}$.)
  \end{corollary}
  \begin{pf}
    Let $Q > 1$. Consider the set
    \[ \{ (x_1, \dots, x_n, y) \in \Z^{n+1}: |x_i-\alpha_iy| \leq Q^{-\nicefrac{1}{n}}, 0 < |y| \leq Q\}.\]
    Let
    \[ \begin{pmatrix}
      L_1(x,y) \\
      \vdots \\
      L_{n+1}(x,y)
    \end{pmatrix} = \underbrace{\begin{pmatrix}
      1 & 0 & \cdots & 0 & -\alpha_1 \\
      0 & \ddots & \ddots & \vdots & \vdots \\
      \vdots & \ddots & \ddots & 0 & -\alpha_{n-2} \\
      \vdots & & \ddots & \ddots & -\alpha_{n} \\
      0 & \cdots & \cdots & 0 & 1\end{pmatrix}}_{\coloneq \widehat{A}}\begin{pmatrix}
        x_1 \\
        \vdots \\
        x_n \\
        y
      \end{pmatrix}.\]
    We use Corollary \ref{cor:1.4.4} (with $A_1 = \dots = A_n = Q^{-\nicefrac{1}{n}}$ and $A_{n+1}=Q$). Then there exists $(x,y) \in \Z^{n+1}\setminus \{ 0\}$ such that
    \begin{align*}
      |x_i-\alpha_iy| \leq Q^{-\nicefrac{1}{n}} \\
      |y| \leq \quad & Q \forall_{1 \leq i \leq n}.
    \end{align*}
    (since $A_1\cdots A_{n+1} = 1 = \det\widehat{A}$). We know that $y \neq 0$ because $y = 0$ would imply $(x,y) = 0$. If $\gcd(x_1, \dots, x_n,y) > 1$, then divide through by the $\gcd$. Also, if $y < 0$, divide throuhg by $(-1)$. Now, for any $Q > 1$, there exists $(x_1, \dots, x_n,y) \in \Zprim{n+1} \setminus \{ 0\}$ such that $0 < y \leq Q$ and $|x_i-\alpha_iy| \leq Q^{-\nicefrac{1}{n}} \forall_{1 \leq i \leq n}$. The assertion now follows through reasoning analogously to Corollary \ref{cor:irr_app}.
  \end{pf}
  Given $\theta \in \R$, we denote
  \[ \|\theta\| = \min\{ |\theta-n|: n \in \Z\}.\]
  From Corollary \ref{cor:1.4.5}, for $\alpha_1,\alpha_2 \in \R$, not both in $\Q$, we get
  \begin{align*}
    & \# \{ y \in \N: \|\alpha_1y\| \leq y^{-\nicefrac{1}{2}}, \|\alpha_2y\| \leq y^{-\nicefrac{1}{2}}\} = \aleph_0 \\
    \implies & \#\{ y \in \N: y\|\alpha_1y\|\|\alpha_2y\| \leq 1\} = \aleph_0
  \end{align*}
  (note that this also holds for $\alpha_1,\alpha_2 \in \Q$).
  \begin{hypothesis}{Littlewood}{littlewood}
    Let $\alpha_1,\alpha_2 \in \R$. For any $\epsilon > 0$, there exists $y \in \N$ such that
    \[ y\|\alpha_1y\|\|\alpha_2y\| < \epsilon.\]
  \end{hypothesis}
  \section{Minkowski's Second Theorem}
  Let $\Lambda$ be a lattice and $C$ a central symmetric convex body in $\R^n$. The \define[Successive Minima of a Lattice with respect to a Central Symmetric Convex Body]{successive minima} $0 < \lambda_1 \leq \lambda_2 \leq \dots \leq \lambda_n$ of $\Lambda$ with respect to $C$ are defined as
  \[ \lambda_i = \min\{ \lambda \in \R_{\geq 0}: \dim\Span_\R(\lambda C \cap \Lambda) \geq i\} \forall_{1 \leq i \leq n}.\]
  \begin{lemma}{Well-definedness and Construction of the Successive Minima}{2.8}
    The successive $0 < \lambda_1 \leq \dots \leq \lambda_n < \infty$ are well-defined. Also, there are $v_1, \dots, v_n \in \Lambda$, which are linearly independent over $\R$, such that $v_i \in \lambda_iC \forall_{1 \leq i \leq n}$.
  \end{lemma}
  \begin{pf}
    Recall that
    \[ \|x\|_C = \min\{ \lambda \in \R_{\geq 0}: x \in \lambda C\} \text{ and } \lambda C = \{ x \in \R^n: \|x\|_C \leq \lambda\}.\]
    Let $\Lambda \setminus \{ 0\} = \{ x_1, \dots\}$, where $0 < \| x_1 \|_C \leq \dots$. Also note that given any $m \in \N$, there are finitely many $i$ such that $m-1 \leq \|x_i\|_C \leq m$ (because $mC$ is bounded and $\lambda$ is discrete). Let $\lambda_1 = \|x_1\|_C$. For $i\in \{ 2, \dots, n\}$, define
    \[ k_i = \min\{ k \in \N: \dim\Span_\R\{ x_1, \dots, x_k\} = i\} \text{, set } v_i=x_i \text{ and let } \lambda_i = \|x_i\|_C.\]
    Clearly, $v_i \in \lambda_iC \cap \Lambda$ and $0 < \lambda_1 \leq \dots \leq \lambda_n < \infty$. Finally, we show that these are indeed the successive minima: By construction, the
    \[ v_1, \dots, v_i \in \lambda_iC \cap \Lambda \text{ are linearly independent over } \R.\]
    Take $0 < \lambda < \lambda_i$ and let $\ell$ be the largest number such that $x_{\ell} \in \lambda C \cap \Lambda$. Then, we have $\ell < k_i$, and by the definition of the $k_i$, it follows that
    \[ \dim\Span_\R(\lambda C \cap \Lambda) = \dim\Span_\R\{ x_1, \dots, x_{\ell}\} < i.\]
    This shows that $\lambda_i$ is indeed the successive minimum.
    \end{pf}
    \begin{theorem}{Minkowski's Second Convex Body Theorem}{mink_sec_conv_bod_thm}
      Let $\Lambda \subseteq \R^n$ be a lattice and $C \subseteq \R^n$ a central symmetric convex body. Further assume that $\lambda_1, \dots, \lambda_n$ are the successive minima of $\Lambda$ with respect to $C$. Then, we have
      \[ \frac{2^n}{n!}\frac{\det\Lambda}{\vol(C)} \leq \lambda_1\cdots\lambda_n \leq 2^n\frac{\det\Lambda}{\vol(C)}.\]
    \end{theorem}
    Before we prove this, let us study a couple of examples to highlight the importance of this result.
    \begin{example}
      Let $n \geq 2$ and
      \[ \mathfrak B = B_1^n(0).\]
      Further let $\Lambda \subseteq \R^n$ be a lattice and $\lambda_1, \dots, \lambda_n$ be the successive minima of $\Lambda$ with respect to $\mathfrak B$. It follows from Lemma \ref{lem:2.8} that there exist $x_1, \dots, x_n \in \Lambda$, which are linearly independent over $\R$ and fulfil $x_i \in \lambda_i\mathfrak B \cap \Lambda$. By construction, we have $\|x_i\|_{\mathfrak B} = \lambda_i \iff |x_i|=\lambda_i$. Theorem \ref{thm:mink_sec_conv_bod_thm} implies
      \[ \prod_{i=1}^n |x_i| \leq \frac{2^n\det\Lambda}{V_n},\]
      where
      \[ V_n = \vol(\mathfrak B) = \begin{cases}
        1 & \text{if } n=1, \\
        \pi & \text{if } n=2, \\
        \frac{2\pi}{n}V_{n-2} & \text{if } n \geq 3.
      \end{cases}\]
      Note that $\{ x_1, \dots, x_n\}$ might not be a basis of $\Lambda$.
    \end{example}
    \begin{example}
      Let $\Lambda = \Z^n$. Take any real $0 < \lambda_1 < \dots < \lambda_n$. Define $C = \{ x \in \R^n: |x_i| \leq \nicefrac{1}{\lambda_i} \forall_{1 \leq i \leq n}\}$. Then $C$ is a central symmetric convex body with
      \[ \vol(C) = \frac{2^n}{\lambda_1\cdots\lambda_n}. \hfill (*)\]
      Let us now show that the $\lambda_j$ are the successive minima of $\Lambda$ with respect to $C$. For $\lambda \geq 0$, we have
      \[ \lambda C = \{ x \in \R^n: |x_i| \leq \nicefrac{\lambda}{\lambda_i} \forall_{1 \leq i \leq n}\}.\]
      Then $e_1, \dots, e_j \in \lambda_jC$ are linearly independent over $\R$. Suppose that $0 < \lambda < \lambda_j$ and $y \in \Z^n \cap \lambda C$. Necessarily, $y_{j}=\dots=y_n = 0$. Therefore,
      \[ \Z^n \cap \lambda C \subseteq \Z^{j-1} \times \{ (0,\dots, 0) \in \Z^{n-j+1}\}.\]
      This cannot contain $j$ linearly independent vectors. Hence, $\lambda_j$ is the successive minimum. Thus,
      \[ \lambda_1\cdots\lambda_n \stackrel{(*)}{=} \frac{2^n}{\vol(C)} = \frac{2^n\det\Lambda}{\vol(C)}.\]
      This proves that the upper bound is sharp.
    \end{example}
    \begin{example}
      Take $\Lambda = \Z^n$, \[C = \left\{ x \in \R^n: \sum_{i=1}^n |x_i| \leq 1\right\}.\]
      \begin{exercise}{}{}
        Show that $\vol(C) = \frac{2^n}{n!}$ and $\lambda_1 = \dots = \lambda_n = 1$.
      \end{exercise}
      This proves that the lower bound in Theorem \ref{thm:mink_sec_conv_bod_thm} is sharp.
    \end{example}
    \newLecture{29th April 2026}
    \begin{lemma}{Convex Hull}{2.10}
    Let $w_1, \dots, w_r \in \R^n$ and define
    \[C(w_1, \dots, w_r) = \left\{ \sum_{i=1}^r x_iw_i: (x_1, \dots, x_r) \in \R^r, \sum_{i=1}^r |x_i|\leq 1\right\}.\]
    Then $C(w_1, \dots, w_r)$ is the smallest convex subset in $\R^n$ which contains $w_1, \dots, w_r$ and is symmetric about $0$.
    We call $C(w_1, \dots, w_r)$ the \define[Convex Hull]{convex hull} of $w_1, \dots, w_r$.
    \end{lemma}
    \begin{pf}
      This is an easy exercise.
    \end{pf}
    \begin{lemma}{Lattice-Base Ordering Compatible with Successive Minima}{2.11}
      Let $\lambda_1, \dots, \lambda_n$ be the successive minima of $\Lambda$ with respect to $C$. Then there exists a basis $\{ b_1, \dots, b_n\}$ of $\Lambda$ such that if
      \[ x = u_1b_1+\dots+u_nb_n\]
      with $\|x\|_C < \lambda_j$ for some $j \in \{ 1, \dots, n\}$, then
      \[ x = u_1b_1+\dots+u_{j-1}b_{j-1}.\]
    \end{lemma}
    \begin{pf}
      Let $\mu_1<\cdots<\mu_s$ be such that $\mu_1 = \lambda_1 = \dots = \lambda_{k_1}$, $\mu_2 = \lambda_{k_1+1} = \dots = \lambda_{k_2}$ and so on with $k_s = n$. \\
      By the definition of the successive minima, $a = 0\forall_{a \in \Lambda: \|a\|_C < \mu_1}$. Since $\mu_2 > \lambda_{k_1}$, by Lemma \ref{lem:2.8}, there exist linearly independent $a_1, \dots, a_n \in \Lambda$ such that $\|a_1\|_C, \dots, \|a_{k_1}\|_C \leq \lambda_{k_1} < \mu_2$. And given any $x \in \Lambda$ with $\|x\|_C < \mu_2$, then
      \[ \exists_{u_1, \dots, u_{k_1} \in \R}: x = u_1a_1 + \dots + u_{k_1}a_{k_1}\]
      (again use the definition of successive minima). We may continue in this manner (extend $a_1, \dots, a_{k_1}$) to obtain linearly independent $a_1, \dots, a_n \in \Lambda$ with the following property:
      \[ \forall_{1 \leq t \leq s}: a_1, \dots, a_{k_{t-1}} \text{ are linearly independent over }\R \text{ and } \|a_j\|_C < \mu_t \forall_{1 \leq j \leq k_{t-1}}.\]
      Given any $x \in \Lambda$ with $\|x\|_C < \mu_t$, we have
      \[ x = u_1a_1 + \dots + u_{k_{t-1}}a_{k_{t-1}}\]
      for some $u_i \in \R$ (!). \\
      Then, by Lemma \ref{lem:lat_base_trans}, we obtain a basis $\{ b_1, \dots, b_n\}$ of $\Lambda$ such that
      \begin{align*}
        b_1 = \quad & \xi_{1,1}a_1 \\
        b_2 = \quad & \xi_{2,1}a_1 + \quad  \xi_{2,2}a_2 \\
        \vdots \qquad \ & \\
        b_n = \quad & \xi_{n,1}a_1 + \quad  \dots + \quad  \xi_{n,n}a_n
      \end{align*}
      with $\xi_{i,j} \in \R, \xi_{i,i} \neq 0 \forall_{1 \leq i,j \leq n}$. This basis satisfies the desired property. Verifying this is an easy exercise.
    \end{pf}
    We shall now prove Minkowski's Second Theorem:
    \begin{pf}[Proof of Theorem \ref{thm:mink_sec_conv_bod_thm}]
      We prove the correctness of the lower and upper bound separately.
      \begin{enumerate}[(i)]
        \item{\textbf{Lower Bound} \\
        Let $v_1, \dots, v_n \in \Lambda$ be linearly independent over $\R$ such that $v_i \in \lambda_i C \forall_{1 \leq i \leq n}$ (Lemma \ref{lem:2.8}). By definition, we have $\lambda_i^{-1}v_i \in C \forall_{1 \leq i \leq n}$. \\
        Define
        \[ D = C(\lambda_1^{-1}v_1, \dots, \lambda_n^{-1}v_n) \subseteq C\]
        and note that this has the properties stated in Lemma \ref{lem:2.10}. Now, let $\phi: \R^n \to \R^n$ be a linear transformation by $\phi(e_i) = \lambda_i^{-1}v_i \forall_{1 \leq i \leq n}$. Then $D = \phi(\mathfrak{B}_1)$, where
        \[ \mathfrak{B}_1 = \left\{ x \in \R^n: \sum_{i=1}^n |x_i| \leq 1\right\}.\]
        Hence,
        \begin{align*}
          \vol(D) = \quad & |\det\phi|\vol(\mathfrak{B}_1) \\
          = \quad & \frac{|\det[v_1\cdots v_n]|}{\lambda_1\cdots\lambda_n}\frac{2^n}{n!} \\
          = \quad & \frac{\det\widetilde{\Lambda}}{\lambda_1\cdots\lambda_n}\frac{2^n}{n!},
        \end{align*}
        where $\widetilde{\Lambda}$ is a lattice with basis $\{ v_1, \dots, v_n\}$. Since $\{ v_1, \dots, v_n\} \subseteq V$, $\widetilde{\Lambda}$ is a sublattice of $\Lambda$.
        This implies that
        \[ \det\widetilde{\Lambda} = [\Lambda:\widetilde{\Lambda}]\det\Lambda \geq \det\Lambda.\]
        We combine everything to get
        \[ \vol(C) \geq \vol(D) = \frac{\det\widetilde{\Lambda}}{\lambda_1\cdots\lambda_n}\frac{2^n}{n!} \geq \frac{\det\Lambda}{\lambda_1\cdots\lambda_n}\frac{2^n}{n!}.\]Thus,
        \[ \lambda_1\cdots\lambda_n \geq \frac{\det\Lambda}{\vol(C)}\frac{2^n}{n!}.\]}
        \item{\textbf{Upper Bound} \\
        Let $\Lambda$ be a lattice inside $\R^n$ and $C$ a central symmetric convex body inside $\R^n$. Given $t > 0$, let
        \[ S(t) = S_F(t) = \{ \eta \in F \mid \exists_{x \in \Lambda}:\eta+x \in tC\}.\]
        where $F$ is the fundamental parallel piped by $\Lambda$.
        The assertion follows from Theorem \ref{thm:volthm_st} by applying it with $t = \nicefrac{\lambda}{2}$. We have
        \begin{align*}
          \det\Lambda = \vol(F) \geq \vol(S(\nicefrac{\lambda}{2})) \geq \frac{\lambda_1}{2}\cdots\frac{\lambda_n}{2}\vol(C),
        \end{align*}
        which can be transformed into the statement.}
      \end{enumerate}
    \end{pf}

  \begin{pf}[Proof of Theorem \ref{thm:2.12}]
      We can reduce the prolem to the case where the basis $\{ b-1, \dots, b_n\}$ given in Lemma \ref{lem:2.11} is $\{ e_1, \dots, e_n\}$ (Exercise sheet 3). In particular, $\Lambda = \Z^n$. Let $1 \leq j \leq n-1$ and $x_1 = (x_{1,1}, x_{1,n}), x_2 = (x_{2,1}, \dots, x_{2,n})$ be two vectors in $tC$ with $t < \nicefrac{\lambda_{j+1}}{2}$, such that $x_1-x_2 \in \Lambda$. Then
      \begin{align*}
        \underbrace{\|x_1-x_2\|_C}_{\in \Z^n} \leq \quad & \|x_1\|_C + \|x_2\|_C \leq 2t < \lambda_{j+1}. \hfill (1)
      \end{align*}
      This implies $x_1-x_2 \in \Z^j \times \{ (0,\dots, 0)\}$. Denote
      \[ S_j(t,z) = \{ (\eta_1, \dots, \eta_n) \in [0,1]^j \mid \exists_{(m_1, \dots, m_j) \in \Z^j}: \|\underbrace{(\eta_1+m_1, \dots, \eta_j+m_j,z)}_{\in tC}\|_C \leq t\}.\]
      Recall that
      \[ S(t) = \{ \eta \in [0,1)^n \mid \exists_{m \in \Z^n}: \|\eta+m\|_C \leq t\}.\]
      Let $\eta \in [0,1)^n$ and $t < \nicefrac{\lambda_{j+1}}{2}$. Supppose that we have $m, m' \in \Z^n$ such that
      \[ \|\eta+m\|_C \leq t \text{ and } \|\eta+m'\|_C \leq t.\]
      Then, by (1), we get $(m_{j+1}, \dots, m_n) = (m'_{j+1}, \dots, m'_n)$. Thus, for each $\eta \in [0,1)^n$ from the definition of $S(t)$, the last $n-1$ coordinates of $m$ are unique. We partition $S(t)$ based on this.
      Write
      \[ S(t) = \bigsqcup_{m_{j+1}, \dots, m_n \in \Z}\bigsqcup_{\eta_{j+1}, \eta_n \in [0,1)} \{ \eta \in [0,1)^n: \exists_{(m_1, \dots, m_j) \in \Z^j: \|\eta+m\|_C \subseteq t}\}.\]
      (Here $m=(m_1, \dots, m_n)$).
      Therefore,
      \begin{align*}
        \vol(S(t)) = \quad & \int_{\R^{n-j}} \vol_J(s_j(t,z)) \dl{z} \forall_{t < \frac{\lambda_{j+1}}{2}}.\hfill (2)
      \end{align*}
    \end{pf}
\newLecture{6th May 2026}

Let \(\Lambda\) be a lattice inside~\(\R^n\) and let \(C\) be a central symmetric convex body.
Given \(t > 0\), let
\begin{equation*}
    S(t)
    = \{ \eta \in F \mid \eta + x \in F \text{ for some } x \in \Lambda\}.
\end{equation*}

\begin{theorem}{Volume of $S(t)$}{volthm_st}
    We have
    \begin{equation*}
        \vol(S(t)) \ge
        \begin{cases*}
            t^n \vol(C) & if \(t \le \frac{\lambda_1}{2}\), \\
            \frac{\lambda_1}{2} \dotsm \frac{\lambda_j}{2} t^{n-j} \vol(C) & if \(\frac{\lambda_j}{2} \le t \le \frac{\lambda_{j+1}}{2}\), \\
            \frac{\lambda_1}{2} \dotsm \frac{\lambda_n}{2} \vol(C) & if \(t \ge \frac{\lambda_n}{2}\).
        \end{cases*}
    \end{equation*}
\end{theorem}

\begin{proof}
    We reduced to the case case where the \glqq nice\grqq basis given in Lem.~2.11 is \(\{e_1, \dots, e_n\}\).
    In particular, \(\Lambda = \Z^n\).
    Let \(1 \le j \le n-1\).
    Denote
    \begin{equation*}
        S_j(t_1 z)
        = \{(\eta_1, \dots, \eta_j) \in [0,1)^j \mid \text{there is } (m_1, \dots, m_j) \in \Z^j \text{ such that } \|(\eta_1 + m_1, \dots, \eta_j + m_j, z)\|_C \le t\}.
    \end{equation*}
    We showed that
    \begin{equation}
        \vol(S(t))
        = \int_{\R^{n-j}} \vol_j(S_j(t, z)) \dl{z} \tag{2}
        \label{eqn:minkowski2:2}
    \end{equation}
    for any \(t < \lambda_{j+1} / 2\).

    Now, take any \(\alpha \ge 1\).
    Claim: \(\vol(S_j(\alpha t, \alpha z)) \ge \vol(S_j(t, z))\) for any \(z \in \R^{n-j}\).

    If the right hand side is~\(0\), then there is nothing to prove.
    Suppose otherwise that there is \(y_0 = (y_{0,1}, \dots, y_{0,j}) \in \R^j\) such that \(\|(y_0, z)\|_C \le t\).
    Given any \(y \in \R^j\) such that \(\|(y, z)\|_C \le t\), then
    \begin{equation}
        \|(y + (\alpha - 1) y_0, \alpha z)\|_C
        = \|(y, z) + (\alpha - 1) (y_0, z)\|_C
        \le \|(y, z)\|_C + (\alpha - 1) \|(y_0, z)\|_C
        \le t + (\alpha - 1) t
        = \alpha t. \tag{3}
        \label{eqn:minkowski2:3}
    \end{equation}
    Let us define a notation: given \(y \in \R^j\), let
    \begin{equation*}
        \{y\}
        = (y_1 - \lf y_1\rf, \dots, y_j - \lf y_j\rf) \in [0,1)^j.
    \end{equation*}
    Then
    \begin{equation*}
        S_j(t, z)
        = \{\{y \}\in [0,1)^j \mid y \in \R^j, \|(y, z)\|_C \le t\}.
    \end{equation*}
    Now, we define
    \begin{equation*}
        W
        = \{\{y + (\alpha - 1) y_0\} \in [0,1)^j \mid y \in \R^j, \|(y, z)\|_C \le t\}.
    \end{equation*}
    Then \(\vol(S_j(t, z)) = \vol(W)\) (check this!).
    Furthermore, by~\cref{eqn:minkowski2:3}, we have
    \begin{equation*}
        W
        \subseteq \{\{y + (\alpha - 1) y_0\} \in [0,1)^n \mid y \in \R^j, \|(y + (\alpha - 1) y_0, \alpha z)\|_C \le \alpha t\}
        = \{\{y\} \in [0,1)^j \mid y \in \R^j, \|(y, \alpha z)\|_C \le \alpha t\}
        = S_j(\alpha t, \alpha z).
    \end{equation*}
    Take the volume to obtain
    \begin{equation*}
        \vol(S_j(t, z))
        = \vol(W)
        \le \vol(S_j(\alpha t, \alpha z)),
    \end{equation*}
    proving the claim.

    Suppose that \(0 < t \le \alpha t < \lambda_{j+1} / 2\).
    Then by \cref{eqn:minkowski2:2} and the claim,
    \begin{equation*}
        \vol(S(\alpha t))
        = \int_{\R^{n-j}} \vol(S_j(\alpha t, z)) \dl{z}
        = \alpha^{n-j} \int_{\R^{n-j}} \vol(S_j(\alpha t, \alpha z)) \dl{z}
        \ge \alpha^{n-j} \int_{\R^{n-j}} \vol(S(t, z)) \dl{z}.
    \end{equation*}
    This gives
    \begin{equation}
        \vol(S(\alpha t)) \\
        \ge \alpha^{n-j} \vol(S(t)). \tag{4}
        \label{eqn:minkowski2:4}
    \end{equation}

    We are ready to prove the theorem.
    When \(t < \lambda_1 / 2\), we have
    \begin{equation*}
        \vol(S(t))
        = \vol(tC)
        = t^n \vol(C).
    \end{equation*}
    This is because if \(\vec{\eta} + \vec{m}, \vec{\eta} + \vec{m}' \in tC\) with \(\vec{m}, \vec{m}' \in \Z^n\), then
    \begin{equation*}
        \|\vec{m} - \vec{m}'\|_C
        = \|(\vec{\eta} + \vec{m}) - (\vec{\eta} + \vec{m}')\|_C
        \le \|\vec{\eta} + \vec{m}\|_C + \|\vec{\eta} + \vec{m}'\|_C
        \le 2t
        < \lambda_1,
    \end{equation*}
    so \(\vec{m} = \vec{m}'\).
    Therefore, \(\vol(S(t)) = \vol(tC) = t^n \vol(C)\) if \(t < \lambda_1 / 2\).
    Since \(S(t) \subseteq S(\lambda_1 / 2)\) for any \(t \le \lambda_1 / 2\), we have \(\vol(S(\lambda_1 / 2)) \ge (\lambda_1 / 2)^n \vol(C)\).

    We now use induction to prove the statement for \(1 \le j \le n-1\).
    Suppose that the statement holds for any \(t \le \lambda_j / 2\) for some~\(j\).
    Consider \(\lambda_j / 2 \le t < \lambda_{j+1} / 2\).
    Set \(t_0 = \lambda_j / 2\) and \(\alpha \in [1, \lambda_{j+1} / \lambda_j)\) such that \(t = \alpha t_0\).
    Then
    \begin{equation*}
        \vol(S(t))
        = \vol(S(\alpha t_0))
        \ge \alpha^{n-j} \vol(S(t_0))
        \ge \alpha^{n-j} \left(\frac{\lambda_1}{2} \dotsm \frac{\lambda_{j-1}}{2}\right) t_0^{n-j+1} \vol(C)
        = t^{n-j} \left(\frac{\lambda_1}{2} \dotsm \frac{\lambda_{j-1}}{2} \frac{\lambda_j}{2}\right) \vol(C).
    \end{equation*}
    And for \(t = \lambda_{j+1} / 2\), this holds by taking the limit since \(S(t) \subseteq S(\lambda_{j+1} / 2)\) if \(t \le \lambda_{j+1} / 2\).
    This completes induction and gives us the statement for all \(t \le \lambda_n / 2\).

    Finally, suppose that \(t \ge \lambda_n / 2\).
    Then \(S(\lambda_n / 2) \subseteq S(t)\), and so
    \begin{equation*}
        \vol(S(t))
        \ge \vol(S(\lambda_n / 2))
        \ge \frac{\lambda_1}{2} \dotsm \frac{\lambda_n}{2} \vol(C).
    \end{equation*}
    (This follows from the case \(t = \lambda_n / 2\)\@.)
\end{proof}

This prove is, largely exaggerated, just a very technical induction.
The proof of Roth's theorem will be also, largely exaggerated, a~very technical induction.

\section{Dual lattice}

Let \(\Lambda\) be a lattice in~\(\R^n\).
We define the dual lattice~\(\Lambda^*\) to be
\begin{equation*}
    \Lambda^*
    = \{\vec{x} \in \R^n \mid \langle\vec{x},y\rangle \in \Z \text{ for all } y \in \Lambda\}.
\end{equation*}
Let \(\{\vec{v}_1, \dots, \vec{v}_n\}\) be a basis of~\(\Lambda\) and let \(A = \begin{pmatrix} \vec{v}_1 & \cdots & \vec{v}_n\end{pmatrix} \).
Then
\begin{equation*}
    \begin{split}
        \Lambda^*
        &= \{\vec{x} \in \R^n \mid \inp{\vec{x}}{y} \in \Z \text{ for all } y \in \Lambda\} \\
        &= \{\vec{x} \in \R^n \mid \inp{\vec{x}}{A z} \in \Z z \in \Z^n\} \\
        &= \{\vec{x} \in \R^n \mid \inp{A^t \vec{x}}{z} \in \Z z \in \Z^n\} \\
        &= \{\vec{x} \in \R^n \mid A^t \vec{x} \in \Z^n\} \\
        &= \{(A^t)^{-1} \vec{w} \mid \vec{w} \in \Z^n\}.
    \end{split}
\end{equation*}
Hence,
\begin{equation*}
    \det \Lambda^*
    = \|\det(A^t)^{-1}\|
    = \frac{1}{\abs{\det A}}
    = \frac{1}{\det \Lambda}
\end{equation*}

\begin{theorem}
    Let \(\Lambda\) be a lattice in~\(\R^n\) and let \(\lambda_1, \dots, \lambda_n\) be its successive minima w.\,r.\,t.\ \(B = \{\vec{x} \in \R^n \mid \|\vec{x}\|_2 \le 1\}\).
    Let \(\Lambda^*\) be its dual, and let \(\lambda_1^*, \dots, \lambda_n^*\) be the successive minima of~\(\Lambda^*\) w.\,r.\,t.~\(B\).
    Then
    \begin{equation*}
        1
        \le \lambda_i \lambda_{n-i+1}^*
        \le c(n)
    \end{equation*}
    for all \(1 \le i \le n\) where \(c(n)\) is a constant that depends only on~\(n\).
\end{theorem}

Successive minima for dual lattices are only considered w.\,r.\,t.\ the unit ball.

\begin{proof}
    We begin with the lower bound.
    Let \(\vec{v}_1, \dots, \vec{v}_n \in \Lambda\) be linearly independent over~\(\R\) and \(\|\vec{v}_i\|_2 = \lambda_i\) for all \(1 \le i \le n\) (from Lem.~2.8).
    Similarly, take \(\vec{v}_1^*, \dots, \vec{v}_n^* \in \Lambda^*\) linearly independent over~\(\R\) with \(\|\vec{v}_i^*\|_2 = \lambda_i^*\) for all \(1 \le i \le n\).
    Take \(1 \le i \le n\).
    Define
    \begin{equation*}
        V_i
        = \{\vec{x} \in \R^n \mid \inp{\vec{x}}{\vec{v}_k} = 0 \text{ for each } 1 \le k \le i\}.
    \end{equation*}
    Since \(\vec{v}_1, \dots, \vec{v}_i\) are linearly independent over~\(\R\), we obtain \(\dim V_i = n-i\).
    Therefore, at least one of \(\vec{v}_1^*, \dots, \vec{v}_{n+1-i}^*\) is not in~\(V_i\).
    Thus, there is \(1 \le i_0 \le i\) and \(1 \le k_0 \le n+1-i\) such that \(\inp{\vec{v}_{i_0}}{\vec{v}_{k_0}^*} \ne 0\).
    From the definition of~\(\Lambda^*\), we have \(\abs{\inp{\vec{v}_{i_0}}{\vec{v}_{i_0}^*}} \in \Z\) and
    \begin{equation*}
        1
        \le \abs{\inp{\vec{v}_{i_0}}{\vec{v}_{i_0}^*}}
        \le \|\vec{v}_{i_0}\|_2 \dotsm \|\vec{v}_{k_0}^*\|_2
        = \lambda_{i_0} \lambda_{k_0}^*
        < \lambda_i \lambda_{n+1-i}^*.
    \end{equation*}

    Next, the upper bound.
    By Theorem \ref{thm:mink_sec_conv_bod_thm}, we have
    \begin{equation*}
        \lambda_1 \dotsm \lambda_n
        \le \frac{2^n}{\vol(B)} \det \Lambda, \qquad
        \lambda_1^* \dotsm \lambda_n^*
        \le \frac{2^n}{\vol(B)} \det \Lambda^*.
    \end{equation*}
    This gives
    \begin{equation*}
        \prod_{i=!}^n \lambda_i \lambda_{n+1-i}^*
        = \frac{4^n}{\vol(B)^2} \det \Lambda \det \Lambda^*
        = \frac{4^n}{\vol(B)^2}.
    \end{equation*}
    Take any \(1 \le i_0 \le n\).
    Then
    \begin{equation*}
        \lambda_{i_0} \lambda_{n+1-i_0}^*
        = \frac{4^n}{\vol(B)^2} \frac{1}{\prod_{i \ne i_0} \underbrace{\lambda_i \lambda_{n+1-i}^*}_{\ge 1}}
        \le \frac{4^n}{\vol(B)^2}. \qedhere
    \end{equation*}
\end{proof}

As an application, we will show that there is no \enquote{small} vector in~\(\Lambda^*\).
This gives a good approximation of vectors in~\(\R^n\) by~\(\Lambda\).

\begin{corollary}
    Let us use the same notation as in Theorem~2.13.
    Suppose that \(0 < R \le \lambda_1^*\).
    Then, given any \(\vec{b} \in \R^n\), there is \(\vec{x} \in \Lambda\) such that
    \begin{equation*}
        \|\vec{b} - \vec{x}\|_2
        \le \frac{n c(n)}{2R}.
    \end{equation*}
\end{corollary}

\begin{proof}
    Since \(\lambda_1^* \ge R\), by Theorem~2.13, we obtain
    \begin{equation*}
        \lambda_n
        \le \frac{c(n)}{\lambda_1^*}
        \le \frac{c(n)}{R}.
    \end{equation*}
    Let \(\vec{v}_1, \dots, \vec{v}_n \in \Lambda\) be linearly independent over~\(\R\) with \(\|\vec{v}_i\|_2 = \lambda_i \le \lambda_n \le c(n) / R\) for all \(1 \le i \le n\).
    Take any \(\vec{b} \in \R^n\), and write \(\vec{b} = \xi_1 \vec{v}_1 + \dots + \xi_n \vec{v}_n\).
    Let \(m_1, \dots, m_n \in \Z\) such that \(\abs{\xi_i - m_i} \le \frac{1}{2}\).
    Then \(\vec{x} = m_1 \vec{v}_1 + \dots + m_n \vec{v}_n \in \Lambda\).
    We have
    \begin{equation*}
        \|\vec{x} - \vec{b}\|_2
        = \|(\xi_1 - m_1) \vec{v}_1 + \dots + (\xi_n - m_n) \vec{v}_n\|_2
        \le \frac{1}{2} (\|\vec{v}_1\|_2 + \dots + \|\vec{v}_n\|_2)
        \le \frac{n}{2} \frac{c(n)}{R}. \qedhere
    \end{equation*}
\end{proof}
\newLecture{8th May 2026}
Recall from last time... \\
Given a lattice $\Lambda$, we defined the dual lattice
\[ \Lambda^* = \{ x \in \R^n: \langle x,y\rangle \in \Z, y \in \Lambda \}.\]
If $\Lambda = A\Z^n$, then $\Lambda^* = (A^t)^{-1}\Z^n$.
\begin{corollary}{}{2.14}
  Using the same notation as in Theorem \ref{thm:2.13}, suppose $0 < R \leq \lambda_1^*$. Then given any $b \in \R^n$, there exists $x \in \Lambda$ such that \[\| b - x \|_2 \leq \frac{nC(n)}{2R}.\]
\end{corollary}
\begin{theorem}{}{2.15}
  Let $\alpha_1, \dots, \alpha_n, \theta_1, \dots, \theta_n \in \R$. Suppose that $\{ 1, \alpha_1, \dots, \alpha_n\}$ are linearly independent over $\Q$. Then, for any constant $\epsilon > 0$, there exist infinitely many $(x_1, \dots, x_n, y) \in \Z^{n+1}$ such that
  \[ |\alpha_iy - x_i-\theta_i| \leq \epsilon \forall_{1 \leq i \leq n}.\]
\end{theorem}
\noindent \textbf{Remark:} Recall Corollary \ref{cor:2.7}, which was a multi-dim Dirichlet's theorem. This is a version with a shift $\theta$ (inhomogeneous), and a nice application of the concept of a dual lattice. Also note that the condition of the linear independence is necessary (exercise).
\begin{pf}
  Let $M$ be a sufficiently large integer. Also let
  \[ A = \begin{pmatrix}
    & & & -\alpha_1 \\
    & I_n & & \cdots \\
    & & & -\alpha_n \\
    0 & \dots & 0 & M^{-1}
  \end{pmatrix}\]
  and
  \[ \Lambda_M = \{ Az: z = (x_1, \dots, x_n,y) \in \Z^{n+1}\} = \{ (x_1-\alpha_1y, \dots, x_n-\alpha_ny, M^{-1}y): (x_1, \dots, x_n, y) \in \Z^{n+1}\}.\]
  Write
  \[ b = \begin{pmatrix}
    \theta_1 \\
    \vdots \\
    \theta_n \\
    2\epsilon
  \end{pmatrix}.\]
  The goal is to show that \[\exists_{z \in \Z^{n+1}}: \|b-Az\| \leq \epsilon. \hfill (*)\]
  This would imply
  \[ \epsilon \geq \sqrt{(x_1-\alpha_1y-\theta_1)^2 + \dots + (x_n-a_ny-\theta_n)^2 + (M^{-1}y-2\epsilon)^2} \geq \begin{cases}
    |x_i-\alpha_iy-\theta_i| & \forall_{1 \leq i \leq n}, \\
    |M^{-1}y-2\epsilon & i=n+1.
  \end{cases}\]
  The latter is equivalent to $\epsilon M \leq y \leq 3\epsilon M$. Hence, we can generate infinitely many $(x_1, \dots, x_n, y) \in \Z^{n+1}$ that satisfy the statement of the theorem. We have
  \[ (A^t)^{-1} = \begin{pmatrix}
    & & & 0 \\
    & I_n & & \vdots \\
    & & & 0 \\
    M\alpha_1 & \dots & M\alpha_n & M
  \end{pmatrix}.\]
  Therefore, our dual lattice is
  \[ \Lambda_M^* = (A^t)^{-1}\Z^{n+1} = \{ (x_1, \dots, x_n, M(\alpha_1x_1+\dots+\alpha_nx_n+y)): x_1, \dots, x_n,y \in \Z\}.\]
  Let
  \[ R = \frac{(n+1)C(n+1)}{2\epsilon}\]
  and
  \[ \mu = \min_{(x_1, \dots, x_n,y) \in S} |\alpha_1x_1+\dots+\alpha_nx_n+y|,\]
  where
  \[ S = \{ (x_1, \dots, x_n,y) \in \Z^{n+1} \setminus \{ 0\}. |x_i| < R \forall_{1 \leq i \leq n}, |\alpha_1x_1+\dots+\alpha_nx_n+y| \leq 1\}.\]
  The set $S$ is clearly finite. Recall that $\alpha_1, \dots, \alpha_n$ are linearly independent over $\Q$. Hence, $\mu > 0$. We choose
  \[ M > \max\left\{ \frac{R}{\mu}, R\right\}.\]
  Then, given any $u \in \Lambda_M^* \setminus \{ 0\}$, we have
  \[ \|u\|_2 = \sqrt{x_1^2+\dots+x_n^2+M^2(\alpha_1x_1+\dots+\alpha_nx_n+y)^2} \geq R,\]
  which we can prove by considering the following three cases:
  \begin{enumerate}[(i)]
    \item $\exists_{1 \leq i \leq n}: |x_i| \geq R \checkmark$.
    \item $\forall_{1 \leq i \leq n}: |x_i| < R, (x_1, \dots, x_n,y) \notin S$ \\
    Use $M > R$.
    \item $\forall_{1 \leq i \leq n}: |x_i| < R, (x_1, \dots, x_n,y) \in S$ \\
    Use $M > \nicefrac{R}{\mu}$.
  \end{enumerate}
  This means that $\lambda_1^* > R$ (first minimum of $\Lambda_M^*$ with respect to the unit ball). We can then apply Corollary \ref{cor:2.14} to deduce that
  \[ \exists_{v = Az \in \Lambda_M^*}: \|b-Az\|_2 \leq \frac{(n+1)C(n+1)}{2R} = \epsilon.\]
  This is the bound $(*)$, which finishes the proof.
\end{pf}
\section{Lagrange's Four Square Theorem}
\begin{theorem}{Lagrange, 1770}{2.16}
  Given $n \in \N$, $\exists_{a,b,c,d \in \Nzero}$ such that
  \[ n = a^2+b^2+c^2+d^2.\]
\end{theorem}
We'll do a Geometry of Numbers-style proof of this. We begin with a simple lemma:
\begin{lemma}{}{2.17}
  Given any $p \in \PP$, there exist $a,b \in \Z$ such that
  \[ a^2+b^2 \equiv -1 \mod p.\]
\end{lemma}
\begin{pf}
  If $p=2$, then $(a,b)=(1,0)$ works. For $p > 2$, define
  \[ \A = \{ a^2 \smod p: 0 \leq a \leq \nicefrac{p}{2}\}.\]
  The elements of $\A$ are distinct because
  \[ a^2 \equiv \widehat{a}^2 \mod p \implies p \mid a^2-\widehat{a}^2 = (a+\widehat{a})(a-\widehat{a}),\]
  from which we can immediately deduce that $p \mid a-\widehat{a}$, i.e. $a \equiv \widehat{a} \mod p$. Therefore,
  \[ \#\A = \frac{p-1}{2}+1 = \frac{p+1}{2}.\]
  Similarly, we define
  \[ \mathfrak{B} = \{ -1-b^2 \smod p: 0 \leq b < \nicefrac{p}{2}\}.\]
  By the same argument,
  \[ \#\mathfrak B = \frac{p+1}{2}.\]
  Therefore, we can use the pigeonhole principle to deduce that $\A \cap \mathfrak B \neq \varnothing$, i.e.
  \[ \exists_{a \in \A, b \in \mathfrak B}: a^2 \equiv -1-b^2 \mod p. \qedhere\]
\end{pf}
We also need the following lemma, the proof of which will be on the next exercise sheet.
\begin{lemma}{}{2.18}
  Let $k_1, \dots, k_n \in \N$, $m \geq 1$. Assume that $L_1(x), \dots, L_m(x) \in \Z[x]$ are linear forms. Define
  \[ \Lambda = \{ x \in \Z^n: L_j(x) \equiv 0 \smod k_j \forall_{1 \leq j \leq m}\}.\]
  Then $\Lambda$ is a lattice with $\det\Lambda \leq k_1\cdots k_m$.
\end{lemma}
\begin{pf}
  Exercise.
\end{pf}
These two lemmata allow us to prove Theorem \ref{thm:2.16}.
\begin{pf}[Proof of Theorem \ref{thm:2.16}]
  For $n=1$, the statement is trivial. Otherwise, write $n = m^2p_1\cdots p_r$, where $p_1, \dots, p_r$. Writing $\widetilde{a} = ma, \dots, \widetilde{d} = md$, it suffices to consider $n=p_1\cdots p_r$. We define
  \[ \{ (u_1, \dots, u_4) \in \Z^4: u_1-a_pu_3-b_pu_4 \equiv 0 \smod p, u_2-b_pu_3+a_pu_4 \equiv 0 \smod p \forall_{\PP \ni p \mid n}\},\]
  where $a_p, b_p$ are such that $a_p^2+b_p^2 \equiv -1 \mod p$ (use Lemma \ref{lem:2.17}). Then $\Lambda$ is a lattice with $\det\Lambda \leq n^2$. Let $\epsilon > 0$ be sufficiently small with respect to $n$. Define
  \[ C = B^4_{\sqrt{2(n-\epsilon)}}(0).\]
  Then
  \[ \vol(C) = \frac{\pi^2}{2}\sqrt{2(n-\epsilon)^4} =2\pi^2(n-\epsilon)^2 >  16n^2 \geq 2^4\det\Lambda.\]
  Now, Minkowski's First Theorem \ref{thm:mink_first_conv_bod_thm} implies that there exists $0 \neq u \in \Lambda \cap C$. In particular,
  \[ 0 < U_1^2+u_2^2+u_3^2+u_4^2 \leq 2(n-\epsilon) <2n.\]
  We also have $u \in \Lambda$, which means that
  \begin{align*} u_1^2+u_2^2+u_3^2+u_4^2 \equiv \quad &  (a_pu_3+b_pu_4)^2+(b_pu_3-a_pu_4)^2+u_3^2+u_4^2 \\
  \equiv \quad & (a_p^2+b_p^2+1)u_3^2 + (a_p^2+b_p^2+1)u_4^2 \\
  \equiv \quad & 0 \hspace{9cm} \mod p\end{align*}
  for each $\PP \ni p \mid n$. \\
  Hence, $p \mid u_1^2+u_2^2+u_3^2+u_4^2 \forall_{\PP \ni p \mid n}$, which implies $n \mid u_1^2+u_2^2+u_3^2+u_4^2$. Combining this with the upper bound on $n$ yields
  \[ n = u_1^2+u_2^2+u_3^2+u_4^2.\]
\end{pf}
\section{Short Vectors in a Lattice}
Apparently, finding the shortest vector in a lattice is an important (and quite difficult) problem in cryptography. Some cryptographic schemes rely on this fact. The secret corresponds to the shortest vector. Though we cannot do this in general, we can find a \glqq pretty short\grqq \ vector (LLL-Algorithm). We will do this next week.
\newLecture{13th May 2026}
\end{document}
